\documentclass[10pt,twocolumn,letterpaper]{article}
\usepackage{cvpr}
\renewcommand{\ttdefault}{cmtt}
%% This file contains a number of tweaks that are typically applied to the main document.
%% They are not enabled by default, but can be enabled by uncommenting the relevant lines.

%%
%% Inline annotations; for predefined colors, refer to "dvipsnames" in the xcolor package:
%% https://tinyurl.com/overleaf-colors
%%
\newcommand{\red}[1]{{\color{red}#1}}
\newcommand{\todo}[1]{{\color{red}#1}}
\newcommand{\TODO}[1]{\textbf{\color{red}[TODO: #1]}}
%%
%% disable for camera ready / submission by uncommenting these lines  
%%
% \renewcommand{\TODO}[1]{}
% \renewcommand{\todo}[1]{#1}

%%
%% work harder in optimizing text layout. Typically shrinks text by 1/6 of page, enable
%% it at the very end of the writing process, when you are just above the page limit
%%
% \usepackage{microtype}

%%
%% fine-tune paragraph spacing
%%
% \renewcommand{\paragraph}[1]{\vspace{.5em}\noindent\textbf{#1.}}

%%
%% globally adjusts space between figure and caption
%%
% \setlength{\abovecaptionskip}{.5em}


%%
%% Allows "the use of \paper to refer to the project name"
%% with automatic management of space at the end of the word
%%
% \usepackage{xspace}
% \newcommand{\paper}{ProjectName\xspace}

%%
%% Commonly used math definitions
%%
% \DeclareMathOperator*{\argmin}{arg\,min}
% \DeclareMathOperator*{\argmax}{arg\,max}

%%
%% Tigthen underline
%%
% \usepackage{soul}
% \setuldepth{foobar}
\usepackage{booktabs}
\definecolor{cvprblue}{rgb}{0.21,0.49,0.74}
\usepackage[breaklinks,colorlinks,allcolors=cvprblue]{hyperref}

\def\paperID{20}
\def\confName{CVPR}
\def\confYear{2026}

\title{What a Camera's Path Knows: Stress-Testing a Pixel-Free Activity Model}

\author{Juan Ignacio Arguello \qquad Valerio Costa \qquad Facundo Lucero\\
Bocconi University, 20600 Deep Learning for Computer Vision}

\begin{document}
\maketitle

\begin{abstract}
CamFormer \cite{xue2026seeing} recognises what a person is doing from the path of their head-mounted camera alone,
without any image. We reproduce its main result on Ego-Exo4D exactly (44.81\% five-way motion-to-text accuracy,
where chance is 20\%) and test four changes. Each is fine-tuned from the released checkpoint, with three training
seeds (one for gaze), and compared with a baseline. (1) Training against wrong answers from the
same recording, through an extra loss term, raises accuracy on procedural activities such as cooking and repair
from 33.4\% to 36.9\%. (2) One model trained on both first-person (Ego-Exo4D) and third-person (DynPose-100K)
trajectories matches the first-person specialist; on third-person video it trails by 3.2 points, a gap our
1,000-question test cannot resolve. (3) Eye gaze as a second input shows no confirmed gain. (4) Per-frame noise of
up to 1\,cm on the camera positions leaves accuracy unchanged, and training with such noise brings no measurable
gain. Code: \url{https://github.com/calerio/dl4cv-seeing-without-pixels}.
\end{abstract}

\section{Introduction}
Picture a pair of smart glasses that records only where it is and which way it faces, thirty times a second, with no
camera image and no audio. From that thin line through a room, can you tell whether the wearer is shooting a
basketball, climbing a wall or fixing a bike?

Xue et al.\ \cite{xue2026seeing} show that you often can. Their model, CamFormer, is a small
Transformer (about 0.6\,M parameters\footnote{598,144 in the released checkpoint; the
paper's Table 2 lists 0.3\,M.}) that reads a camera trajectory and maps it into the same space as a frozen
CLIP text encoder \cite{radford2021clip}, so that a trajectory can be matched against written descriptions of what
happens. On Ego-Exo4D \cite{grauman2024egoexo4d} it picks the right description out of five 44.8\% of the time.

The authors are frank about where the idea strains (their supplementary 3.3). Step-by-step \emph{procedural}
activities such as cooking or bike repair score far lower than sports, because the camera motion is subtle and local. And
first-person and third-person cameras get separate models. We took those limits as our starting points and asked
how far a pixel-free model can be pushed. Our contributions:
\begin{itemize}
\item An exact reproduction of the main result from the released code and checkpoints (Sec.~\ref{sec:repro}).
\item \textbf{Harder wrong answers} in training, drawn from the same recording, which improve procedural accuracy
by 3.5 points (Sec.~\ref{sec:hn}).
\item \textbf{One model for both viewpoints}, which matches the first-person specialist but leaves the
third-person comparison unresolved (Sec.~\ref{sec:egoexo}).
\item \textbf{Eye gaze} as a second input, with no confirmed gain (Sec.~\ref{sec:gaze}).
\item \textbf{Noise on the camera positions}, which leaves accuracy unchanged up to 1\,cm of noise per
coordinate (Sec.~\ref{sec:noise}).
\end{itemize}
We wrote down each comparison's rule before reading its result, with two exceptions. In the hard-negatives
experiment the rule was written after we had seen the first seed's test score, and the design itself followed an
earlier run's test score (Sec.~\ref{sec:hn}). In the two-viewpoint experiment the corrected rerun was decided after
the first round's results were known, although the rule it is judged by was fixed before any run
(Sec.~\ref{sec:egoexo}). All numbers come from code in our repository.

\section{Related Work}
\label{sec:related}
\paragraph{Activity from motion, without pixels.} Head and camera motion has long been used to recognise
first-person activity; Kitani et al.\ \cite{kitani2011fast} cluster the egomotion of sports videos into actions
without labels. IMU2CLIP \cite{moon2022imu2clip} trains a motion-sensor encoder against frozen CLIP encoders so that
sensor signals can be matched with text. CamFormer \cite{xue2026seeing} applies the same contrastive idea to
camera trajectories, and is the model we study.

\paragraph{Hard negatives.} InfoNCE \cite{oord2018cpc} and CLIP \cite{radford2021clip} use the rest of the batch as
wrong answers. Robinson et al.\ \cite{robinson2021hard} show that negatives similar to the anchor give better
representations. For egocentric video, EgoNCE \cite{lin2022egovlp} adds clips from the same video and close in time
as extra negatives, and treats clips that share a verb and a noun as positives. Our within-recording loss follows
the same intuition for trajectories, as a separate loss term, and our verb rule removes same-verb pairs from the negatives, a
milder version of EgoNCE's positives.

\paragraph{First and third person.} Ego-Exo \cite{li2021egoexo} transfers video representations from third-person
to first-person video, and Ego-Exo4D \cite{grauman2024egoexo4d} films the same activities from both views. We ask
whether one model can serve both views for camera trajectories, where the two views differ in sensor, scale and frame rate.

\paragraph{Gaze and robustness.} Li et al.\ \cite{li2018eye} model gaze and actions together in first-person video
and report that gaze helps recognition. For robustness we follow Hendrycks and Dietterich
\cite{hendrycks2019corruptions}, who average accuracy over several corruption levels.

\section{Methodology}
\label{sec:method}
\paragraph{CamFormer.} The model turns each camera pose into a token (position relative to the clip's middle frame
plus a rotation, with gravity direction known for Aria), adds a position embedding and passes the sequence through
a 4-layer Transformer (width 128, 4 heads). The output for the labelled window is a single vector $m$. The caption
goes through a frozen CLIP text encoder, giving $t$. Both are normalised to unit length.

\paragraph{Batch loss.} For a batch of $N$ clip--caption pairs, the authors train with the symmetric InfoNCE loss,
with similarity $s_{ij} = m_i^\top t_j / \tau$ and $\tau = 0.07$:
\begin{equation}
\textstyle\mathcal{L}_{\text{batch}} = -\frac{1}{2N}\sum_{i=1}^{N}\left[\log\frac{e^{s_{ii}}}{\sum_{j\in\mathcal{N}_i} e^{s_{ij}}}
+ \log\frac{e^{s_{ii}}}{\sum_{j\in\mathcal{N}_i} e^{s_{ji}}}\right].
\label{eq:batch}
\end{equation}
In the authors' setup $\mathcal{N}_i$ is the whole batch. With our verb rule (setup C and later), $\mathcal{N}_i$
leaves out every other pair whose caption has the same main verb as caption $i$, so that two clips of ``cuts the
onion'' are not pushed apart.

\paragraph{Within-recording loss.} Our extra term (Sec.~\ref{sec:hn}) has the same form, but each pair competes
only with pairs from its own recording $r(i)$ whose main verb $v$ differs:
\begin{equation}
\mathcal{N}^{\text{take}}_i = \{i\} \cup \{\, j : r(j) = r(i),\ v(j) \neq v(i) \,\},
\end{equation}
and pairs with no such competitor are left out. The training loss is
$\mathcal{L} = \mathcal{L}_{\text{batch}} + \lambda\,\mathcal{L}_{\text{take}}$ with $\lambda = 1$.

\paragraph{Extra inputs.} Two experiments add an input to the token sequence, both through layers that start at
zero, so that the released model is unchanged at the start of training. For the joint model (Sec.~\ref{sec:egoexo})
the domain token is one learned vector per camera type (ego or exo), added to every token. For gaze
(Sec.~\ref{sec:gaze}) a linear layer maps four numbers per frame (gaze yaw, pitch, log depth and a validity flag)
to the token width, and the result is added to that frame's token.

\paragraph{Position noise.} For Sec.~\ref{sec:noise}, each frame's position $p_t$ is replaced by
$p_t + \epsilon_t$ with $\epsilon_t \sim \mathcal{N}(0, \sigma^2 I_3)$, independently per frame, before the
relative positions are computed; rotations are untouched.

\section{Datasets}
\label{sec:data}
No video enters any of our models: we use only the camera trajectories, captions and eye gaze that the authors and
Ego-Exo4D release.

\paragraph{Ego-Exo4D (first person).} People wearing Aria glasses do one of eight activities: basketball, soccer, rock climbing, dancing, playing music, cooking, bike repair
and health procedures (a COVID rapid test or CPR). The glasses' own tracking gives metric camera poses, which
the authors resample to 20 per second. Each training example pairs a window of up to 8\,s of trajectory with one
narration of what the wearer does at that moment, 159,186 pairs in all. The test has 7,079 five-way questions from
671 recordings (2,621 procedural, 4,458 physical), built from the validation recordings; for this reason we never
use validation to choose anything. For the gaze experiment we use Ego-Exo4D's eye-gaze recordings.

\paragraph{DynPose-100K (third person).} Ordinary internet videos whose camera poses were estimated from the
images, here with ViPE, so their scale is only approximate and there is no gravity direction. After the authors'
filter on pose scale, 88,151 clip--caption pairs remain for training; the test has 1,000 five-way questions.

\paragraph{Nymeria (zero-shot only).} Four sets of 1,000 questions about body posture, hands and arms, legs and
feet, and focus of attention, which we use only to check whether a change transfers to data the model was not
trained on.

\paragraph{Task.} Each test question shows one trajectory clip and five captions; the model must pick the caption
that describes the clip (motion-to-text). The wrong captions come from the same recording but describe a different
action, which makes the test hard: the setting, the person and the style of movement are all shared. Accuracy is
reported overall and split into \emph{physical} (sports, dance, music) and \emph{procedural} (cooking, repair, health)
activities.

\begin{table}[t]
\centering\small
\setlength{\tabcolsep}{3pt}
\begin{tabular}{lccc}
\toprule
 & Poses from & Train pairs & Test qs. \\
\midrule
Ego-Exo4D & Aria glasses & 159,186 & 7,079 \\
DynPose-100K & ViPE, from video & 88,151 & 1,000 \\
Nymeria & Aria glasses & -- & $4\times1{,}000$ \\
\bottomrule
\end{tabular}
\caption{Datasets. Aria poses are metric; ViPE poses are estimated from video. All test questions are five-way; chance is 20\%.}
\label{tab:data}
\end{table}

\section{Experiments and Results}
\label{sec:exp}
\paragraph{Protocol.} Every experiment fine-tunes the authors' released 8-second checkpoint for 20 epochs with
their settings (learning rate $10^{-4}$, batch 1,024), always against a baseline that receives the same extra
training without our change, so that a gain is not just ``more training''. Each setup is trained with three seeds, except the gaze runs (one seed, Sec.~\ref{sec:gaze}).
Differences are tested with a paired two-level bootstrap that resamples training seeds and then test recordings
(not questions, since questions from one recording are correlated), 10,000 times, and we report 95\% intervals.
With only three seeds this interval is rough, so we also quote per-seed numbers where a claim rests on them.

\subsection{Reproduction}
\label{sec:repro}
From the released features the Ego-Exo4D test gives 44.81\% motion-to-text, matching the paper's 44.8. Running
the released checkpoint end to end on our cluster gives the same 44.81\% (procedural 33.61\%). The DynPose-100K
figure with ViPE poses (46.3) and all four zero-shot Nymeria figures also match. One number does not: DynPose with
the dataset's own poses gives 39.2 from the released features, against 36.2 in the paper. We re-ran the 44.81
check before every new batch of training, as a guard against silent changes to the code.

\subsection{Same-recording wrong answers}
\label{sec:hn}
Training and test disagree about what a wrong answer looks like. At test time, the wrong captions come from the \emph{same} recording:
same person, same kitchen, different step. In training, as is standard for contrastive models, the wrong answers
are simply the other clips in the batch, drawn from the whole dataset. With a batch of 1,024, only about one of
the 1,023 wrong answers comes from the same recording; we call such same-recording wrong answers \emph{hard negatives}. The model practises telling cooking from basketball, then
gets tested on telling ``C places the jar of salt in the kitchen cabinet'' from ``C wipes the skillet with the kitchen towel'' (a real test question; C is the camera wearer).

We close the gap in two steps. First, batches are filled with groups of eight clips from the same recording,
at least 4\,s apart, so that hard negatives are present (the spacing holds within a group; two groups from one recording can land in the same batch). Second, an extra loss term (weight 1) contrasts each clip only
against the other clips of its own recording. In both losses, a clip whose caption has the same main verb is not used as a wrong answer, so that a repeated
action is not punished (Sec.~\ref{sec:method}). We call this setup B$'$, and
compare it with the authors' setup (A), with the verb rule alone (C), and with the grouping but no extra loss (G).

\begin{table}[t]
\centering\small
\begin{tabular}{lccc}
\toprule
Setup & Procedural & Physical & Overall \\
\midrule
Released model & 33.61 & 51.39 & 44.81 \\
A (authors' setup) & 33.36 & 51.87 & 45.01 \\
C (verb rule only) & 33.41 & 52.01 & 45.12 \\
G (grouping only) & 33.73 & 51.90 & 45.17 \\
B$'$ (grouping + loss) & \textbf{36.91} & \textbf{53.39} & \textbf{47.29} \\
\bottomrule
\end{tabular}
\caption{Same-recording negatives, mean of three seeds, motion-to-text accuracy (\%).}
\label{tab:hn}
\end{table}

B$'$ beats the baseline on procedural questions by +3.55 points (95\% interval +2.37 to +4.83) and wins in every
seed (37.09, 36.44, 37.20 against 33.27, 33.61, 33.19), Table~\ref{tab:hn}. It also gains on physical activities
(+1.53, +0.31 to +2.64), so the procedural gain does not come at their expense. The verb rule alone has no measurable
effect (C$-$A overall +0.11, $-$0.57 to +0.74).

The run with grouping but no extra loss (G) shows where the gain comes from. G is indistinguishable from C on procedural questions
(+0.32, $-$0.60 to +1.23), while B$'$ beats G by +3.18 (+1.95 to +4.49). Putting hard negatives in the batch is not enough on its own, because in the standard loss they are a handful among
a thousand; we measured that they carry under 3\% of the weight on negatives. The extra term needs the grouping to
have pairs to work on, and the grouping does little without the term.

The gain shows in the motion-to-text direction only: in text-to-motion (a caption picks among
five trajectories) B$'-$A is +0.74 on procedural with an interval that includes zero, although both loss terms are
symmetric; we did not find out why. And it does not transfer zero-shot to Nymeria, where every interval includes
zero. B$'$ trains on the same construction the test uses, so part of the gain may be specific to this test format.
The design also followed a test result: the move from grouping alone to B$'$ came after grouping alone scored
44.84 on the test (an earlier, unseeded run with a different negative mask, not G).

\subsection{One model for both viewpoints}
\label{sec:egoexo}
CamFormer comes as two models: one for first-person (ego) trajectories from Aria glasses, one for third-person
(exo) trajectories estimated from ordinary videos (DynPose-100K). We trained a joint model
on both datasets without (U0) and with (U1) a domain token telling it which kind of camera it is reading, and
compared it with an ego specialist (S-ego), an exo specialist (S-exo), and E$\to$X, the ego model fine-tuned on exo
data only, which shares the joint model's inputs and starting weights.

A first round of runs fed every exo clip as if it were filmed at 30 frames per second; judging from the number of
poses per clip, about a third are at 24 or 25. We fixed this and retrained all five setups; the numbers below are from the corrected round. None of the
version-to-version differences was resolved.

\begin{table}[t]
\centering\small
\begin{tabular}{lcc}
\toprule
Setup & Ego test & Exo test \\
\midrule
S-ego (ego specialist) & 44.99 & -- \\
S-exo (exo specialist) & -- & 46.53 \\
E$\to$X (exo only, from ego) & 32.58 & 46.37 \\
U0 (joint, no token) & 45.29 & 43.87 \\
U1 (joint, domain token) & 45.02 & 43.33 \\
\bottomrule
\end{tabular}
\caption{One model for two viewpoints, mean of three seeds, motion-to-text (\%). Ego: Ego-Exo4D, 7,079 questions;
exo: DynPose-100K, 1,000 questions.}
\label{tab:egoexo}
\end{table}

On first-person video the joint model matches the specialist: U1$-$S-ego is +0.03 ($-$0.61 to +0.68), well inside the
$\pm$1.5 margin we set in advance. On third-person video it is 3.2 points behind the specialist ($-$6.49 to +0.20).
The interval includes zero, but it also includes a gap of more than six points, so we cannot claim the joint model
matches. With only 1,000 exo questions and seed-to-seed swings of up to four points (U1: 41.7 to 45.8), this test
cannot resolve differences of that size. The domain token adds nothing measurable (U1$-$U0 on exo $-$0.53,
$-$3.61 to +2.78).

The largest gap in the table is a side effect. After training on exo data alone, E$\to$X drops from about 45 to 32.6 on
the ego test: it forgets the first-person skill it started with. The joint models, which also see ego data, stay at about 45.

\subsection{Eye gaze as a second input}
\label{sec:gaze}
Aria glasses also record where the eyes look. This looked promising for procedural tasks, where the head may stay still
while the eyes move from the knife to the onion. We added gaze as a second input in two ways, keeping the released
encoder frozen and training only a gaze layer (GF), or training everything (GT), each with a control that receives
gaze from a different moment or person (``shuffled'').

A pilot gate on training data, fixed in advance, was not met, so only one seed was trained and no claim is made.
On procedural questions all four single-seed comparisons include zero (GT against no gaze: +0.72, $-$0.39 to
+1.84). The point estimates are positive, so the question stays open, but this is no evidence that gaze helps. We
also could not verify the alignment of gaze and pose in time, so the fair summary is ``no gain with gaze as
aligned here''.

\subsection{Noise on the camera positions}
\label{sec:noise}
Camera tracking always has some error. We asked two questions: how much accuracy CamFormer loses when the
positions it reads are slightly wrong, and whether training on noisy positions reduces that loss. We added independent Gaussian noise to the $x$, $y$ and $z$ of every frame, at standard
deviations of 0.5, 1, 2 and 5\,cm per axis, and compared the authors' setup (A) with a model trained with 0.5\,cm
noise on half its clips (N). Following the standard protocol for corruption robustness
\cite{hendrycks2019corruptions}, the main comparison is accuracy averaged over all four levels, fixed before the
runs.

\paragraph{Why 0.5\,cm, and what it is not.} We chose the training level to match the accuracy Meta reports for
Aria trajectories, a global error of no more than 1.5\,cm in room-scale scenes \cite{engel2023aria}: 0.5\,cm on
each axis is about 0.9\,cm in 3D. We were aware that this matches the \emph{size} of real error but not its
\emph{shape}. Aria's figure is an error of whole trajectories, and its uncorrected error is described as drift
that accumulates \cite{engel2023aria}, which is smooth, whereas our noise is independent on every frame. Per-frame
jitter is one of the most common augmentations for motion sensors \cite{wu2025moda}, but it has been found not to
help a wrist-sensor classifier \cite{um2017augmentation} and to hurt pose estimation from body-worn sensors
\cite{wu2025moda}. We chose it knowingly, as the simplest common perturbation, and treat the result as a stress
test of the model rather than a simulation of real tracking failure. N is also tested on the same kind of noise it was trained on, which
the corruption benchmark avoids \cite{hendrycks2019corruptions}.

\begin{table}[t]
\centering\small
\setlength{\tabcolsep}{4pt}
\begin{tabular}{lccccc}
\toprule
 & clean & 0.5\,cm & 1\,cm & 2\,cm & 5\,cm \\
\midrule
A & 44.93 & 44.80 & 44.91 & 44.28 & 39.71 \\
N & 44.89 & 44.91 & 44.95 & 44.70 & 39.89 \\
\bottomrule
\end{tabular}
\caption{Accuracy (\%) under position noise per axis, overall, mean of three seeds and three noise draws.}
\label{tab:noise}
\end{table}

On the second question no gain shows: averaged over the four levels, N is ahead of A by only +0.21 ($-$0.08 to +0.51), and
no cost on clean inputs is resolved either ($-$0.03, $-$0.44 to +0.38). The interval is narrow, so any gain above about half
a point would probably have shown.

On the first question (Table~\ref{tab:noise}), without any special training, accuracy
does not move up to 1\,cm of noise, and slips by less than a point at 2\,cm, even though the noise is
large next to the motion. In a sample of 20 training recordings the median distance the camera moves between frames ranges from
0.2\,cm (a person standing at a kitchen counter) to 3.7\,cm (basketball); at 1\,cm per axis the noise on each
frame-to-frame step is about 2.4\,cm, larger than the median step in 13 of the 20. In the slowest takes the noise
is many times the motion, yet accuracy holds. Only at 5\,cm does accuracy fall, by about five points.

Our reading, which this experiment does not test directly, is that CamFormer relies less on the fine motion
between frames than on the larger shape of the path over seconds, such as the head crossing a room or bending down to a counter. Independent jitter averages out over a clip; walking across a room does not. Smooth, drifting
noise would attack exactly that shape, and is the natural next test.

\section{Conclusions}
A camera's path alone says a lot about what its wearer is doing, and the model that reads it held up better than we
expected. The one change that clearly helped was to train the way the test asks: against wrong answers from the
same recording, with a loss that makes them count. Merging viewpoints works for first-person video but leaves an
unresolved question on third-person video, gaze shows no confirmed gain, and centimetre-level jitter barely moves
the accuracy. If the next test, smooth drifting noise, hurts the model where per-frame jitter did not, that would confirm that it
reads the shape of the path over seconds rather than its frame-to-frame detail.

{\small
\bibliographystyle{ieeenat_fullname}
\bibliography{refs}
}
\end{document}
